\section{Conclusion}
\label{sec-conclude}

An ontology store is asked for patterns, for hierarchies, and for meaning.
This paper described \sys, which answers all three inside one disk-based
engine.
A B+ tree store of three permutations, with crabbing insert and delete, a range
scan that holds no latch between rows, and a bulk loader that fills pages,
answers selective patterns faster than an LSM store, a Java B+ tree store, and a
relational store, and is within a factor of two of the best on large results.
Labelling the class DAG with interval sets turns ``everything under $c$'' into
a few range scans of the same tree, at 2.0 intervals per class on the Gene
Ontology and a fraction of the cost of materializing the closure, and makes
membership in a class's subtree a binary search, which a filtered vector search
needs.
A sorted index orders and does not say what is close in meaning, which the
bandwidth argument and the measurements agree on; exact pre-filtering through
the intervals is the right answer for classes up to about a thousand members
and a masked graph search for the largest.
The whole engine runs in a browser from a 0.85\,MB module.

\stitle{Limitations}.
\sys implements a \sparql subset without \texttt{OPTIONAL}, \texttt{UNION},
property paths, or aggregates, and a reasoner only for
\texttt{rdfs:subClassOf}; a user who needs the full language needs one of the
other systems we measured.
It is the slowest to load, single-threaded, and hungry for memory while
loading, because its optimizer's exact statistics sit in hash tables.
Decoding an id costs about 0.7\,$\mu$s and dominates large results, which the
decode cache hides on a warm store and does not remove.
Labels are rebuilt, not maintained, when the hierarchy changes.
The vectors come from one ontology and one small model that embeds names only,
the instances of the Gene Ontology are synthetic, vectors are held in memory,
and the semantic measurements are single-threaded.
Queries are terms of the ontology; free text would need an embedding model that
\sys does not link.
We did not find a vector index that makes the hierarchy order unnecessary:
neighbors by meaning and descendants are different sets (a third of a term's
ten nearest neighbors are within two \texttt{is\_a} steps, two thirds are
not), which is why the two are combined at query time and not at index time.
